\makeheading{Week 3}
\section{Permutation and Randomization Tests}
\begin{itemize}
    \item All the previous tests have made some kind of distributional assumption
          for the response measurements, such as $ Y_{ij} \sim \N{\mu_j,\sigma^2} $ or
          $ Y_{ij}\sim \Binomial{1,\pi_j} $.
    \item It would be preferable to have a test that does not rely on \emph{any}
          assumptions.
    \item This is precisely the purpose of permutation and randomization tests.
          \begin{itemize}
              \item These tests are \emph{non-parametric} and rely on resampling.
              \item The motivation is that if $ \mathbf{H}_0 $: $ \theta_1=\theta_2 $
                    is true, any random rearrangement of the data is \emph{equally likely
                        to have been observed}. If $ \mathbf{H}_0 $ is true, then we have a single population/distribution
                    from which our data has been drawn.
              \item With $ n_1 $ and $ n_2 $ units in each condition, there are
                    \[ \Binomial{n_1+n_2}{n_1}=\Binomial{n_1+n_2}{n_2} \]
                    arrangements of the $ n_1+n_2 $ observations into two groups of size $ n_1 $
                    and $ n_2 $ respectively.
                    \[ n_1=n_2=50\implies\Binomial{n_1+n_2}{n_1}=\Binomial{100}{50}=1.0089\times 10^{29} \]
          \end{itemize}
    \item A true \textbf{permutation test} considers \emph{all possible rearrangements}
          of the original data.
          \begin{itemize}
              \item The test statistic $ t $ is calculated on the original data and on every one of
                    its rearrangements.
              \item This collection of test statistic values generate the empirical null distribution.
          \end{itemize}
    \item In a \textbf{randomization test}, we do not
          consider all possible rearrangements.
          \begin{itemize}
              \item We just consider a large number $ N $ of them.
              \item We use this in practice instead of a permutation test because
                    the exact permutation tests have too many permutations to consider.
          \end{itemize}
\end{itemize}
\begin{framed}
    \textbf{Randomization Test Algorithm}
    \begin{enumerate}
        \item Collect response observations in each condition.
              \[ \Set{y_{11},y_{21},\ldots,y_{n_1 1}}\to \hat{\theta}_1 \]
              \[ \Set{y_{12},y_{22},\ldots,y_{n_2 2}}\to \hat{\theta}_2 \]
        \item Calculate the test statistic $ t $ on the original data.
              \[ t=\hat{\theta}_1-\hat{\theta}_2\quad\text{or}\quad t=\frac{\hat{\theta}_1}{\hat{\theta}_2} \]
        \item Pool all the observations together and randomly sample (without replacement)
              $ n_1 $ observations which will be assigned to ``Condition 1'' and the remaining $ n_2 $
              observations that are assigned to ``Condition 2.'' Repeat this $ N $ times.
              \[ \Set{y_{11}^\star,y_{21}^\star,\ldots,y_{n_1 1}^\star}\to \hat{\theta}_1^\star \]
              \[ \Set{y_{12}^\star,y_{22}^\star,\ldots,y_{n_2 2}^\star}\to \hat{\theta}_2^\star \]
        \item Calculate the test statistic $ t_k^\star $ on each of the ``shuffled'' datasets, $ k=1,2,\ldots,N $.
              \[ t_k^\star=\hat{\theta}_{1,k}^\star-\hat{\theta}_{2,k}^\star\quad\text{or}\quad t_k^\star=\frac{\hat{\theta}_{1,k}^\star}{\hat{\theta}_{2,k}^\star}  \]
        \item Compare to $ t $ to $ \Set*{t_1^\star,t_2^\star,\ldots,t_N^\star} $, the empirical
              null distribution, and calculate the $ p $-value:
              \[ p\text{-value}=\frac{\text{\# of $t^\star$'s that are at least as extreme as $t$}}{N}=\frac{1}{N} \sum_{k=1}^{N} \Ind*{t_k^\star\text{ at least as extreme as $t$}}  \]
              \begin{itemize}
                  \item $ \mathbf{H}_0 $: $ \theta_1=\theta_2 $ versus $ \mathbf{H}_\text{A} $: $ \theta_1\ne \theta_2 $. If $ t=\hat{\theta}_1-\hat{\theta}_2 $, then
                        the $ p $-value is:
                        \[ p\text{-value}=\frac{1}{N} \sum_{k=1}^{N} \Ind*{t_k^\star\ge \abs{t}\cup t_k^\star\le -\abs{t}} \]
                  \item $ \mathbf{H}_0 $: $ \theta_1\ge \theta_2 $ versus $ \mathbf{H}_\text{A} $: $ \theta_1<\theta_2 $. If $ t=\hat{\theta}_1-\hat{\theta}_2 $, then
                        the $ p $-value is:
                        \[ p\text{-value}=\frac{1}{N} \sum_{k=1}^{N} \Ind*{t_k^\star\le t} \]
                  \item $ \mathbf{H}_0 $: $ \theta_1\le \theta_2 $ versus $ \mathbf{H}_\text{A} $: $ \theta_1>\theta_2 $. If $ t=\hat{\theta}_1-\hat{\theta}_2 $, then
                        the $ p $-value is:
                        \[ p\text{-value}=\frac{1}{N} \sum_{k=1}^{N} \Ind*{t_k^\star\ge t} \]
              \end{itemize}
    \end{enumerate}
\end{framed}
\begin{Example}{Pokémon Go}{}
    \begin{itemize}
        \item Suppose that Niantic Inc, is experimenting with two different promotions within Pokémon
              Go:
              \begin{itemize}
                  \item Condition 1: Give users nothing.
                  \item Condition 2: Give users 200 free Pokécoins.
                  \item Condition 3: Give users a 50\% discount on Shop purchases.
              \end{itemize}
        \item In a small pilot experiment, we randomize $ n_1=n_2=n_3=100 $ users to each condition.
        \item For each user, we record the amount of real money (in USD) they spend in the 30 days
              following the experiment.
        \item The data summaries are:
              \begin{itemize}
                  \item $ \bar{y}_1=\$10.74 $, $ Q_{y_1}(0.5)=\$9 $.
                  \item $ \bar{y}_2=\$9.53 $, $ Q_{y_2}(0.5)=\$8 $.
                  \item $ \bar{y}_3=\$13.41 $, $ Q_{y_3}(0.5)=\$10 $.
              \end{itemize}
              Using R, we performed a randomization test with $ N=10\,000 $ with respect to the mean
              we found that the control and free coin conditions did not significantly differ.
              But there was a significant increase in the amount of money spent in the discount condition
              relative to the other two.

              \vspace{2mm}

              The hypotheses that we tested to determine these conclusions were:
              \begin{tightcenter}
                  $ \mathbf{H}_0 $: $ \mu_1=\mu_2 $ versus $ \mathbf{H}_\text{A} $: $ \mu_1\ne \mu_2 $\\
                  $ \mathbf{H}_0 $: $ \mu_1\ge \mu_2 $ versus $ \mathbf{H}_\text{A} $: $ \mu_1<\mu_2 $
              \end{tightcenter}
              Interestingly, when you run these same tests, but on the basis of the median, we find
              no significant difference between any of the conditions.
        \item \href{https://github.com/Hextical/university-notes/blob/master/year-3/semester-3/STAT 430/code/W3/Randomization_test.R}{[R Code] \texttt{Randomization\_test}}
    \end{itemize}
\end{Example}
\chapter{Experiments with More than Two Conditions}
\section*{Anatomy of an ``A/B/$ m $'' Test}
\begin{itemize}
    \item One design factor at $ m $ levels.
    \item We will now consider a design and analysis of an experiment consisting
          of more than two experimental conditions --- or what many data scientists
          broadly refer to as ``A/B/$m$ Testing.''
          \begin{itemize}
              \item Canonical A/B/$ m $ test:
                    \begin{figure}[!htbp]
                        \centering
                        \begin{tikzpicture}[x=0.75pt,y=0.75pt,yscale=-1.5,xscale=1.5]
                            %Rounded Rect [id:dp8976617722511913] 
                            \draw  [draw opacity=0][fill={rgb, 255:red, 208; green, 2; blue, 27 }  ,fill opacity=1 ] (10,19.53) .. controls (10,15.12) and (13.58,11.53) .. (18,11.53) -- (81.3,11.53) .. controls (85.72,11.53) and (89.3,15.12) .. (89.3,19.53) -- (89.3,43.53) .. controls (89.3,47.95) and (85.72,51.53) .. (81.3,51.53) -- (18,51.53) .. controls (13.58,51.53) and (10,47.95) .. (10,43.53) -- cycle ;
                            %Rounded Rect [id:dp05576705853009267] 
                            \draw  [draw opacity=0][fill={rgb, 255:red, 74; green, 144; blue, 226 }  ,fill opacity=1 ] (102,19.07) .. controls (102,14.65) and (105.58,11.07) .. (110,11.07) -- (173.3,11.07) .. controls (177.72,11.07) and (181.3,14.65) .. (181.3,19.07) -- (181.3,43.07) .. controls (181.3,47.48) and (177.72,51.07) .. (173.3,51.07) -- (110,51.07) .. controls (105.58,51.07) and (102,47.48) .. (102,43.07) -- cycle ;
                            %Rounded Rect [id:dp1316893319581255] 
                            \draw  [draw opacity=0][fill={rgb, 255:red, 65; green, 117; blue, 5 }  ,fill opacity=1 ] (191,18.33) .. controls (191,13.92) and (194.58,10.33) .. (199,10.33) -- (262.3,10.33) .. controls (266.72,10.33) and (270.3,13.92) .. (270.3,18.33) -- (270.3,42.33) .. controls (270.3,46.75) and (266.72,50.33) .. (262.3,50.33) -- (199,50.33) .. controls (194.58,50.33) and (191,46.75) .. (191,42.33) -- cycle ;
                            %Rounded Rect [id:dp1823408740718886] 
                            \draw  [draw opacity=0][fill={rgb, 255:red, 144; green, 19; blue, 254 }  ,fill opacity=1 ] (281,18.33) .. controls (281,13.92) and (284.58,10.33) .. (289,10.33) -- (352.3,10.33) .. controls (356.72,10.33) and (360.3,13.92) .. (360.3,18.33) -- (360.3,42.33) .. controls (360.3,46.75) and (356.72,50.33) .. (352.3,50.33) -- (289,50.33) .. controls (284.58,50.33) and (281,46.75) .. (281,42.33) -- cycle ;

                            % Text Node
                            \draw (49.65,31.53) node  [color={rgb, 255:red, 255; green, 255; blue, 255 }  ,opacity=1 ] [align=left] {CLICK ME};
                            % Text Node
                            \draw (141.65,31.07) node  [color={rgb, 255:red, 255; green, 255; blue, 255 }  ,opacity=1 ] [align=left] {CLICK ME};
                            % Text Node
                            \draw (230.65,30.33) node  [color={rgb, 255:red, 255; green, 255; blue, 255 }  ,opacity=1 ] [align=left] {CLICK ME};
                            % Text Node
                            \draw (320.65,30.33) node  [color={rgb, 255:red, 255; green, 255; blue, 255 }  ,opacity=1 ] [align=left] {CLICK ME};
                        \end{tikzpicture}
                        \caption{Canonical Button Colour Test.}
                    \end{figure}

                    What colour maximizes click-through rate?
          \end{itemize}
    \item Other, more tangible, examples:
          \begin{itemize}
              \item Netflix.
              \item Etsy.
          \end{itemize}
    \item Typically, the goal of such an experiment is to decide which condition is
          optimal with respect to some metric of interest $ \theta $. This could be a:
          \begin{itemize}
              \item mean
              \item proportion
              \item variance
              \item quantile
              \item technically any statistic that can be calculated from sample data
          \end{itemize}
    \item From a design standpoint, such an experiment is \emph{very} similar to
          a two-condition experiment.
          \begin{enumerate}
              \item Choose a metric of interest $ \theta $ which addresses the question
                    you are trying to answer.
              \item Determine the response variable $ y $ that must be measured on each
                    unit to estimate $ \hat{\theta} $.
              \item Choose the design factor $ x $ and the $ m $ levels you will experiment with.
              \item Choose $ n_1,n_2,\ldots,n_m $ and assign units to conditions at random.
              \item Collect the data and estimate the metric of interest in each condition:
                    \[ \hat{\theta}_1,\hat{\theta}_2,\ldots,\hat{\theta}_m \]
          \end{enumerate}
    \item Determining which condition is optimal typically involves a series of pairwise
          comparisons: $ t $-tests, $ z $-tests, or randomization tests.
    \item But it is useful to begin such an investigation with a \emph{gatekeeper}
          test (test of overall equality) which serves to determine whether
          there is \emph{any} difference between the $ m $ experimental conditions.
          Formally, we phrase such a question as the following statistical hypothesis:
          \begin{tightcenter}
              $ \mathbf{H}_0 $: $ \theta_1=\theta_2=\cdots=\theta_m $ versus $ \mathbf{H}_\text{A} $: $ \theta_j\ne \theta_k $ for some $ j\ne k $
          \end{tightcenter}
          In the case of means:
          \begin{tightcenter}
              $ \mathbf{H}_0 $: $ \mu_1=\mu_2=\cdots=\mu_m $ versus $ \mathbf{H}_\text{A} $: $ \mu_j\ne \mu_k $ for some $ j\ne k $
          \end{tightcenter}
          In the case of proportions:
          \begin{tightcenter}
              $ \mathbf{H}_0 $: $ \pi_1=\pi_2=\cdots=\pi_m $ versus $ \mathbf{H}_\text{A} $: $ \pi_j\ne \pi_k $ for some $ j\ne k $
          \end{tightcenter}
\end{itemize}
\section{Comparing Means in Multiple Conditions}
\begin{itemize}
    \item We assume that our response variable follows a normal distribution,
          and we assume that the mean of the distribution depends on the condition
          in which we take the measurements, and that the variance
          is the same across all conditions.
          \[ Y_{ij}\sim \N{\mu_j,\sigma^2}\quad\text{for }i=1,2,\ldots,n_j\text{ and }j=1,2,\ldots,m \]
    \item We use an $ F $-test to test for means:
          \begin{tightcenter}
              $ \mathbf{H}_0 $: $ \mu_1=\mu_2=\cdots=\mu_m $ versus $\mathbf{H}_\text{A}$: $ \mu_j\ne \mu_k $ for some $ j\ne k $
          \end{tightcenter}
\end{itemize}
\subsection{The \texorpdfstring{$ F $}{F}-test for Overall Significance in a Linear Regression}
\begin{itemize}
    \item In particular, we use the $ F $-test for overall significance in an
          \emph{appropriately defined linear regression model}:
          \begin{itemize}
              \item The \emph{appropriately defined linear regression model}
                    in this situation is one in which the response
                    variable depends on $ m-1 $ indicator variables:
                    \[ x_{ij}=\begin{cases*}
                            1 & if unit $ i $ is in condition $ j $ \\
                            0 & otherwise
                        \end{cases*}\quad\text{for }j=1,2,\ldots,m-1. \]
              \item For a particular unit $ i $, we adopt the model:
                    \[ Y_{i}=\beta_0+\beta_1x_{i1}+\beta_2x_{i2}+\cdots+\beta_{m-1}x_{i,m-1}+\varepsilon_i \]
                    \begin{itemize}
                        \item $ Y_i = $ response observation for unit $ i=1,2,\ldots,N=\sum_{j=1}^{m} n_j $.
                        \item $ \varepsilon_i = $ random error term which we assume follows a $ \N{0,\sigma^2} $ distribution
                              independently for all $ i=1,2,\ldots,N $.
                        \item Because we're about to do a regression analysis, the usual \emph{residual diagnostics} are relevant.
                    \end{itemize}
              \item In this model the $ \beta $'s are unknown parameters,
                    and we interpret them in the context of the
                    following expectations:
                    \begin{itemize}
                        \item Expected response in condition $ m $:
                              \[ \E{Y_i\given x_{i1}=x_{i2}=\cdots=x_{i,m-1}=0}=\beta_0=\mu_m \]
                        \item Expected response in condition $ j $:
                              \[ \E{Y_i\given x_{ij}=1}=\beta_0+\beta_j=\mu_j\quad\text{for }j=1,2,\ldots,m-1 \]
                    \end{itemize}
                    \begin{itemize}
                        \item $ \beta_0 $ is the expected response in condition $ m $.
                        \item $ \beta_j $ is the expected difference in response value in condition $ j $
                              versus condition $ m $ for $ j=1,2,\ldots,m-1 $.
                              \begin{align*}
                                  \mu_1     & =\beta_0+\beta_1      \\
                                  \mu_2     & =\beta_0+\beta_2      \\
                                            & \vdotswithin{=}       \\
                                  \mu_{m-1} & = \beta_0+\beta_{m-1} \\
                                  \mu_m     & =\beta_0
                              \end{align*}
                    \end{itemize}
              \item Based on these assumptions $ \mathbf{H}_0 $: $ \theta_1=\theta_2=\cdots=\theta_m $
                    is true if and only if $ \beta_1=\beta_2=\cdots=\beta_{m-1}=0 $, and hence is equivalent
                    to testing:
                    \begin{tightcenter}
                        $ \mathbf{H}_0 $: $ \beta_1=\beta_2=\cdots=\beta_{m-1}=0 $ versus $ \mathbf{H}_\text{A} $: $ \beta_j\ne 0 $ for some $ j$
                    \end{tightcenter}
              \item This hypothesis corresponds, as noted, to the $ F $-test for overall significance in the model.
          \end{itemize}
    \item In regression parlance, the test statistic is the ratio of the
          regression mean squares (MSR)
          to the mean squared error (MSE) in a standard regression-based analysis of variance (ANOVA):
          \[ t=\frac{\text{MSR}}{\text{MSE}} \]
    \item In our setting we can more intuitively think of the test
          statistic as comparing the response variability
          between conditions to the response variability within conditions:
          \begin{itemize}
              \item Average response in condition $ j $: $ \displaystyle \bar{y}_{\bullet j}=\frac{1}{n_j} \sum_{i=1}^{n_j} y_{ij} $.
              \item Overall average response: $ \displaystyle \bar{y}_{\bullet\bullet}=\frac{1}{N} \sum_{j=1}^{m} \sum_{i=1}^{n_j} y_{ij}=\frac{1}{N} \sum_{j=1}^{m} n_j\bar{y}_{\bullet j} $.
              \item Quantifies variability \underline{\emph{between}} conditions:
                    $ \displaystyle \SSC=\sum_{j=1}^{m} \sum_{i=1}^{n_j} \bigl(\bar{y}_{\bullet j}-\bar{y}_{\bullet\bullet}\bigr)^2 $.
              \item Quantifies variability \underline{\emph{within}} conditions:
                    $ \displaystyle \SSE=\sum_{j=1}^{m} \sum_{i=1}^{n_j} \bigl(y_{ij}-\bar{y}_{\bullet j}\bigr)^2 $.
              \item Quantifies \underline{\emph{overall}} variability:
                    $ \displaystyle \SST=\sum_{j=1}^{m} \sum_{i=1}^{n_j} \bigl(y_{ij}-\bar{y}_{\bullet\bullet}\bigr)^2=\SSC+\SSE $.
          \end{itemize}
    \item The null distribution for this test is $ F(m-1,N-m) $.
    \item $ p\text{-value}=\Prob{T\ge t} $ where $ T \sim F(m-1,N-m) $.
    \item If $ \mathbf{H}_0 $: $ \mu_1=\cdots=\mu_m $ is true, then $ \E*{\MSC}=\sigma^2 $ and $ \E*{\MSE}=\sigma^2 $.
\end{itemize}
\begin{table}[!htbp]
    \centering
    \caption{ANOVA Table}
    \begin{NiceTabular}{|l|c|c|c|c|}
        \toprule
        Source    & SS                     & d.f.    & MS                                                               & Test Statistic                                                   \\
        \midrule
        Condition & $ \SSC $ & $ m-1 $ & $ \MSC=\SSC/(m-1) $ & \multirow{2}{*}{$ t=\MSC/\MSE $} \\
        Error     & $ \SSE $ & $ N-m $ & $ \MSE=\SSE/(N-m) $                                                               \\
        \midrule
        Total     & $ \SST $ & $ N-1 $\\
        \bottomrule
    \end{NiceTabular}
\end{table}
\subsection{Example: Candy Crush Boosters}
\begin{Example}{Candy Crush Boosters}{}
    \begin{itemize}
        \item Candy Crush is experimenting with three different versions of
              in-game ``boosters'': the lollipop hammer, the jelly fish, and the colour bomb.
        \item We randomize each user to one of these three conditions ($ n_1 = 121 $,
              $ n_2 = 135 $, $ n_3 = 117$) and they
              receive (for free) 5
              boosters corresponding to their condition.
              Interest lies in evaluating the effect
              of these different boosters on the length of time a user plays the game.
        \item Let $ \mu_j $ represent the average length of game play
              (in minutes) associated with booster condition
              $ j = 1, 2, 3 $. While interest lies in finding the condition
              associated with the longest average length
              of game play, here we first rule out the possibility
              that booster type does not influence the length
              of game play (i.e., $ \mu_1=\mu_2=\mu_3 $).
        \item In order to do this we fit the linear regression model:
              \[ Y=\beta_0+\beta_1 x_1+\beta_2 x_2+\varepsilon \]
              where $ x_1 $ and $ x_2 $ are indicator variables indicating
              whether we observe a particular value of the response in the
              jelly fish or colour bomb conditions, respectively.
              The lollipop hammer is therefore the reference condition.
        \item In R, we found that the test statistic for testing:
              \begin{tightcenter}
                  $ \mathbf{H}_0 $: $ \mu_1=\mu_2=\mu_3 $ versus $ \mathbf{H}_\text{A} $: $ \mu_j\ne \mu_k $ for some $ j\ne k $
              \end{tightcenter}
              was $ t=851.8947 $ and the null distribution was $ T \sim F(2,370) $. The corresponding
              $ p $-value was:
              \[ p\text{-value}=\Prob{T\ge 851.8947}=3.28\times 10^{-139} \]
        \item Therefore, we have very strong evidence against $ \mathbf{H}_0 $ and conclude that the average
              length of game play is not the same in the three booster conditions.
        \item \href{https://github.com/Hextical/university-notes/blob/master/year-3/semester-3/STAT 430/code/W3/Comparing_multiple_means.R}{[R Code] \texttt{Comparing\_multiple\_means}}
    \end{itemize}
\end{Example}
\section{Comparing Proportions in Multiple Conditions}
\begin{itemize}
    \item As is always the case when comparing proportions is of interest, we assume that our response variable
          is binary:
          \[ Y_{ij}=\begin{cases*}
                  1 & if unit $ i $ in condition $ j $ performs an action of interest         \\
                  0 & if unit $ i $ in condition $ j $ does not perform an action of interest
              \end{cases*}i=1,2,\ldots,n_j;j=1,2,\ldots,m \]
    \item $ Y_{ij}\sim \Binomial{1,\pi_j} $
          where $ \pi_j $ is the probability of a unit in condition $ j $ performing the action.
    \item We use a \textbf{chi-squared test of independence} (Pearson $ \chi^2 $ test) to test for proportions:
          \begin{tightcenter}
              $ \mathbf{H}_0 $: $ \pi_1=\pi_2=\cdots=\pi_m $ versus $ \mathbf{H}_\text{A} $: $ \pi_j\ne \pi_k $ for some $ j\ne k $.
          \end{tightcenter}
\end{itemize}
\subsection{The Chi-squared Test of Independence}
\begin{itemize}
    \item We use the chi-squared test of independence as a test for `no association'
          between two categorical variables that are summarized in a \emph{contingency table}.
    \item We apply this methodology here to test the independence of the binary outcome
          (whether a unit performs the action of interest) and the particular condition they are in.
    \item To start, let's assume that $ m=2 $, and let's use the \hyperref[ex:optimizely_ex1]{Optimizely experiment} as a reference.
          \begin{itemize}
              \item If $ \pi_1=\pi_2=\pi $, then we would expect the conversion rate in each condition to be the same.
              \item An estimate of the pooled conversion rate in this case is $ \hat{\pi}=679/17514=0.0388 $.
              \item Let $ X= $ number of conversions in a condition with $ n $ units, therefore $ X \sim \Binomial{n,\pi} $
                    where $ \E{X}=n\pi $.
              \item Therefore, we would expect $ n_1\hat{\pi}=8872(0.0388)=343.96 $ conversions in condition 1,
                    and $ n_2\hat{\pi}=8642(0.0388)=335.04 $ conversions in condition 2.
              \item The chi-squared test formally evaluates if the difference between what
                    was observed and what is expected under the null hypothesis is large enough
                    to be considered \emph{significantly} different.
              \item The \emph{general} $ 2\times 2 $ contingency table for a scenario like this
                    is shown in~\Cref{general_22_contingency}.
                    \begin{table}[!htbp]
                        \centering
                        \caption{A General $ 2\times 2 $ Contingency Table}\label{general_22_contingency}
                        \begin{NiceTabular}{cc|cc|c}
                            \multicolumn{2}{c}{}        & \multicolumn{2}{c}{\emph{Condition}} &                                                                                       \\
                            \multicolumn{2}{c}{}        & 1                             & 2                         & \multicolumn{1}{c}{}                                      \\
                            \cmidrule{2-5}
                            \multirow{2}{*}{\emph{Conversion}} & Yes                           & $O_{1,1}$                 & $O_{1,2}$                 & $O_1$                         \\
                            & No                            & $O_{0,1}$                 & $O_{0,2}$                 & $O_0$                         \\
                            \cmidrule{2-5}
                            \multicolumn{1}{c}{}        & \multicolumn{1}{c}{}          & \multicolumn{1}{c}{$n_1$} & \multicolumn{1}{c}{$n_2$} & \multicolumn{1}{c}{$n_1+n_2$}
                        \end{NiceTabular}
                    \end{table}
                    \begin{itemize}
                        \item $ O_{\ell,j} $: observed number of conversions ($ \ell=1 $),
                              and the observed number of non-conversions ($ \ell=0 $) in condition $ j=1,2 $.
                        \item $ O_\ell $: overall number of conversions ($ \ell=1 $) or non-conversions
                              ($ \ell=0 $)
                    \end{itemize}
              \item So,
                    \[ \hat{\pi}=\frac{O_1}{n_1+n_2}\quad \text{ and }\quad 1-\hat{\pi}=\frac{O_0}{n_1+n_2} \]
                    represent the overall proportions of units that did or did not convert, and they
                    are estimates of overall conversion and non-conversion rates.
              \item Let $ E_{1,j} $ and $ E_{0,j} $
                    represent the expected number of conversions and non-conversions in condition
                    $ j=1,2 $,
                    \[ E_{1,j}=n_j\hat{\pi}\quad \text{ and }\quad E_{0,j}=n_j(1-\hat{\pi}) \]
                    \begin{itemize}
                        \item This is what we expect if $ \mathbf{H}_0 $: $ \pi_1=\pi_2 $ is true.
                    \end{itemize}
              \item The $ \chi^2 $ test statistic compares the observed count
                    in each cell to the corresponding expected
                    count, and is defined as
                    \[ T=\sum_{\ell=0}^{1} \sum_{j=1}^{2} \frac{(O_{\ell_j}-E_{\ell,j})^2}{E_{\ell,j}}\sim \chi^2(1)  \]
              \item $ p\text{-value} = \Prob{T\ge t} $ where $ T \sim \chi^2(1) $.
              \item Returning to the Optimizely example, the \emph{expected} table is~\Cref{optimizely_contingency}.
                    \begin{table}[!htbp]
                        \centering
                        \caption{$ 2\times 2 $ Contingency Table for Optimizely's Homepage Experiment}\label{optimizely_contingency}
                        \begin{NiceTabular}{cc|cc|c}
                            \multicolumn{2}{c}{}        & \multicolumn{2}{c}{\emph{Condition}} &                                                                                       \\
                            \multicolumn{2}{c}{}        & 1                                      & 2                          & \multicolumn{1}{c}{}                                     \\
                            \cmidrule{2-5}
                            \multirow{2}{*}{\emph{Conversion}} & Yes                                    & $343.96$                   & $335.04$                   & $679$                       \\
                            & No                                     & $8528.04$                  & $8306.96$                  & $16835$                     \\
                            \cmidrule{2-5}
                            \multicolumn{1}{c}{}        & \multicolumn{1}{c}{}                   & \multicolumn{1}{c}{$8872$} & \multicolumn{1}{c}{$8642$} & \multicolumn{1}{c}{$17514$}
                        \end{NiceTabular}
                    \end{table}
              \item And the resultant test statistic and $p$-value are:
                    \[ t=\frac{(280-343.96)^2}{343.96} +\frac{(399-335.04)^2}{335.04}+\frac{(8592-8528.04)^2}{8528.04}+\frac{(8243-8306.96)^2}{8306.96}=25.075    \]
                    \[ p\text{-value}=\Prob{T\ge 25.075}=5.52\times 10^{-7}\quad\text{where }T \sim \chi^2(1) \]
          \end{itemize}
    \item Let's now extend this for $ m>2 $.
          \begin{itemize}
              \item We've used the chi-squared test is a test of `no association' between the binary outcome (whether
                    a unit performs the action of interest) and the particular condition they are in.
                    \begin{itemize}
                        \item But there is no requirement that there be only two conditions.
                        \item Here we generalize the test to any number of experimental conditions.
                    \end{itemize}
              \item The information associated with this test can be summarized in a $ 2\times m $ contingency table as seen in~\Cref{general_2m_contingency}.
                    \begin{table}[!htbp]
                        \centering
                        \caption{A General $ 2\times m $ Contingency Table}\label{general_2m_contingency}
                        \begin{NiceTabular}{cc|cccc|c}
                            \multicolumn{2}{c}{}        & \multicolumn{4}{c}{\emph{Condition}} &                                                                                                                                           \\
                            \multicolumn{2}{c}{}        & 1                                      & 2                         & $\cdots$                  & $m$      & \multicolumn{1}{c}{}                                                   \\
                            \cmidrule{2-7}
                            \multirow{2}{*}{\emph{Conversion}} & Yes                                    & $O_{1,1}$                 & $O_{1,2}$                 & $\cdots$ & $O_{1,m}$                 & $O_1$                                      \\
                            & No                                     & $O_{0,1}$                 & $O_{0,2}$                 & $\cdots$ & $O_{0,m}$                 & $O_0$                                      \\
                            \cmidrule{2-7}
                            \multicolumn{1}{c}{}        & \multicolumn{1}{c}{}                   & \multicolumn{1}{c}{$n_1$} & \multicolumn{1}{c}{$n_2$} & $\cdots$ & \multicolumn{1}{c}{$n_m$} & \multicolumn{1}{c}{$N=\sum_{j=1}^{m} n_j$}
                        \end{NiceTabular}
                    \end{table}
                    \begin{itemize}
                        \item \# of conversions ($ \ell=1 $) or non-conversions ($ \ell=0 $) is condition
                              $ j=1,2 $.
                        \item $ \hat{\pi}=O_1/N $.
                        \item $ 1-\hat{\pi}=O_0/N $.
                    \end{itemize}
              \item We compare each of the observed frequencies $ O_{1,j} $ with the corresponding expected
                    frequency $ E_{\ell,j} $.
                    \[ E_{1,j}=n_j\hat{\pi}\quad\text{and}\quad E_{0,j}=n_j(1-\hat{\pi}) \]
                    \begin{itemize}
                        \item Expected number of conversions/non-conversions in condition $ j $
                              assuming $ \mathbf{H}_0 $: $ \pi_1=\pi_2=\cdots=\pi_m $ is true.
                    \end{itemize}
              \item The $ \chi^2 $ test statistic
                    compares the observed count in each cell to the corresponding expected
                    count, and is defined as:
                    \[ T=\sum_{\ell=0}^{1} \sum_{j=1}^{m} \frac{(O_{\ell,j}-E_{\ell,j})^2}{E_{\ell_j}}\sim \chi^2(m-1)  \]
              \item $ p\text{-value} = \Prob{T\ge t} $ where $ T \sim \chi^2(m-1) $.
          \end{itemize}
\end{itemize}
\subsection{Example: Nike SB Ads}
\begin{Example}{Nike SB Ads}{}
    \begin{itemize}
        \item Suppose that Nike is running an ad campaign for Nike SB, their
              skateboarding division, and the
              campaign involves $ m = 5 $ different
              video ads that are being shown in Facebook newsfeeds.
        \item A video ad is `viewed' if it is watched for longer than 3 seconds,
              and interest lies in determining
              which ad is most popular and hence most profitable
              by comparing the viewing rates of the five
              different videos.
        \item We show each of these 5 videos to $ n_1=5014 $, $ n_2=4971 $, $ n_3=5030 $,
              $ n_4=5007 $, and $ n_5=4980 $ users, and summarize the results in~\Cref{nike_contingency}.
              \begin{center}
                  \captionsetup{type=table}
                  \captionof{table}{A $ 2\times 5 $ Observed Contingency Table for the Nike Example}\label{nike_contingency}
                  \begin{NiceTabular}{cc|ccccc|c}
                      \multicolumn{2}{c}{}  & \multicolumn{5}{c}{\emph{Condition}} &                                                                                                                                                                              \\
                      \multicolumn{2}{c}{}  & 1                                      & 2                          & 3                          & 4                          & 5                          & \multicolumn{1}{c}{}                                     \\
                      \cmidrule{2-8}
                      \multirow{2}{*}{\emph{View}} & Yes                                    & $160$                      & $95$                       & $141$                      & $293$                      & $197$                      & $886$                       \\
                      & No                                     & $4854$                     & $4876$                     & $4889$                     & $4714$                     & $4783$                     & $24116$                     \\
                      \cmidrule{2-8}
                      \multicolumn{1}{c}{}  & \multicolumn{1}{c}{}                   & \multicolumn{1}{c}{$5014$} & \multicolumn{1}{c}{$4971$} & \multicolumn{1}{c}{$5030$} & \multicolumn{1}{c}{$5007$} & \multicolumn{1}{c}{$4980$} & \multicolumn{1}{c}{$25002$}
                  \end{NiceTabular}
              \end{center}
        \item The overall watch rate (and its complement) are:
              \[ \hat{\pi}=\frac{O_1}{N} =\frac{886}{25002}=0.0354\quad\text{and}\quad 1-\hat{\pi}=\frac{24116}{25002} =0.9649 \]
        \item We multiply $ n_j $ by $ \hat{\pi} $ and $ (1-\hat{\pi}) $ for $ j=1,2,3,4,5 $
              to get the expected cell frequencies in~\Cref{nike_expected}.
              \begin{center}
                  \captionsetup{type=table}
                  \captionof{table}{A $ 2\times 5 $ Expected Contingency Table for the Nike Example}\label{nike_expected}
                  \begin{NiceTabular}{cc|ccccc|c}
                      \multicolumn{2}{c}{}  & \multicolumn{5}{c}{\emph{Condition}} &                                                                                                                                                                              \\
                      \multicolumn{2}{c}{}  & 1                                      & 2                          & 3                          & 4                          & 5                          & \multicolumn{1}{c}{}                                     \\
                      \cmidrule{2-8}
                      \multirow{2}{*}{\emph{View}} & Yes                                    & $177.68$                   & $176.16$                   & $178.25$                   & $177.43$                   & $176.48$                   & $886$                       \\
                      & No                                     & $4836.32$                  & $4794.84$                  & $4851.75$                  & $4829.57$                  & $4803.52$                  & $24116$                     \\
                      \cmidrule{2-8}
                      \multicolumn{1}{c}{}  & \multicolumn{1}{c}{}                   & \multicolumn{1}{c}{$5014$} & \multicolumn{1}{c}{$4971$} & \multicolumn{1}{c}{$5030$} & \multicolumn{1}{c}{$5007$} & \multicolumn{1}{c}{$4980$} & \multicolumn{1}{c}{$25002$}
                  \end{NiceTabular}
              \end{center}
        \item The resultant test statistic and $ p $-value (where $ T \sim \chi^2(4) $) are:
              \[ t=\sum_{\ell=0}^{1} \sum_{j=1}^{m} \frac{(O_{\ell,j}-E_{\ell,j})^2}{E_{\ell,j}} =129.1686 \]
              \[ p\text{-value}=\Prob{T\ge 129.1686}=5.86\times 10^{-27} \]
        \item Therefore, we reject $ \mathbf{H}_0 $: $ \pi_1=\pi_2=\cdots=\pi_5 $ and conclude that
              the ``watch-rate'' is not the same for each of the video ads.
        \item \href{https://github.com/Hextical/university-notes/blob/master/year-3/semester-3/STAT 430/code/W3/Comparing_multiple_proportions.R}{[R Code] \texttt{Comparing\_multiple\_proportions}}
    \end{itemize}
\end{Example}
